\documentclass[11pt,a4paper]{article}

%% PACHETE MATEMATICE
\usepackage{amsmath,amssymb,amsfonts}
\usepackage{amsthm}
\usepackage{mathpazo}
\usepackage{eulervm}
\usepackage{enumerate}
\usepackage{color}
\usepackage{graphicx}
\usepackage[all]{xy}
\usepackage{xcolor}
\usepackage{framed}
\usepackage[unicode]{hyperref}
\setcounter{MaxMatrixCols}{20}
\usepackage{multicol}
%\usepackage[romanian]{babel}


%% ASPECT
\linespread{1.05}
\usepackage[T1]{fontenc}
\usepackage[utf8x]{inputenc}
\usepackage[margin=2cm]{geometry}
% \usepackage[color]{showkeys}
\usepackage{anyfontsize}
\usepackage{titlesec}
\titleformat*{\section}{\large\bfseries}

\usepackage{fancyhdr}
\pagestyle{fancy}
\rhead{\small \color{gray} Notițe de Adrian Manea}
\lhead{\small \color{gray} Aero-BCD, 2017---2018, M. Olteanu \& L. Costache}


\usepackage[autostyle = false, style = german]{csquotes}
\usepackage[romanian]{babel}
\newcommand\qq{\enquote}

%% COMENZILE MELE
\renewcommand*{\proofname}{Dem.:}
\newcommand\bb{\mathbb}
\newcommand\dr{\mathrm}
\newcommand\kal{\mathcal}
\newcommand\fk{\mathfrak}
\newcommand\op{\oplus}
\newcommand\ot{\otimes}
\newcommand\ds{\displaystyle}
\newcommand\id{\indent}
\newcommand\nid{\noindent}
\newcommand\seq{\subseteq}
\newcommand\sneq{\subsetneq}
\newcommand\speq{\supseteq}
\newcommand\spneq{\supsetneq}
\newcommand\rar{\rightarrow}
\newcommand\Rar{\Rightarrow}
\newcommand\xrar{\xrightarrow}
\newcommand\lar{\leftarrow}
\newcommand\Lar{\Leftarrow}
\newcommand\xlar{\xleftarrow}
\newcommand\lrar{\leftrightarrow}
\newcommand\Lrar{\Leftrightarrow}
\newcommand\ideal{\trianglelefteq}
\newcommand\ai{\text{ a.\^{i}. }}
\newcommand\sk{(\textit{Schi\c{t}\u{a}}) }
\newcommand\rhup{\rightharpoonup}
\newcommand\rhdn{\rightharpoondown}
\newcommand\lhup{\leftharpoonup}
\newcommand\lhdn{\leftharpoondown}
\newcommand\vid{\emptyset}
\newcommand\wh{\widehat}
\newcommand\ol{\overline}
\newcommand\NN{\mathbb{N}}
\newcommand\RR{\mathbb{R}}
\newcommand\ZZ{\mathbb{Z}}
\newcommand\QQ{\mathbb{Q}}
\newcommand\CC{\mathbb{C}}

%% STILURILE MELE
\newtheoremstyle{dotless}{}{}{\itshape}{}{\bfseries}{:}{ }{}

\theoremstyle{dotless}
\newtheorem{theorem}{Teorem\u{a}}[section]
\newtheorem{proposition}{Propozi\c{t}ie}[section]
\newtheorem{lemma}{Lem\u{a}}[section]
\newtheorem{corollary}{Corolar}[section]
\newtheorem{exercise}{Exerci\c{t}iu}

\newtheoremstyle{dotlessdr}{}{}{}{}{\bfseries}{:}{ }{}
\theoremstyle{dotlessdr}
\newtheorem{example}{Exemplu}[section]
\newtheorem{definition}{Defini\c{t}ie}[section]
\newtheorem{remark}{Observa\c{t}ie}[section]




\begin{document}

\begin{center}
  {\large \textbf{Seminar 1 \\ Recapitulare șiruri}}
\end{center}

\vspace{1cm}

\section{Topologia dreptei reale}

\id\id Analiza matematică și studiul numerelor reale se face și cu ajutorul
reprezentării geometrice a acesteia, anume \emph{axa numerelor reale}. Prin
definiție, axa este o dreaptă pe care s-a fixat \emph{o origine, o unitate %
și un sens pozitiv}. Numerele reale se reprezintă pe axă corespunzător cu
valoarea lor, corespondența fiind dată de unitatea fixată. De aceea, putem
observa că pe mulțimea numerelor reale (și pe axă, implicit), avem o
\emph{relație de ordine}, care face ca mulțimea numerelor reale să devină
chiar \emph{total ordonată}.

Prin definiție, spunem că numărul real $a$ este mai mic decît numărul real
$b$, notat $a < b$, dacă diferența $a - b$ este mai mică decît zero, punctul
distins stabilit la fixarea originii.

Mai mult, putem relaxa această definiție și să folosim inegalitatea nestrictă,
$a \leq b$. Așadar, perechea $(\RR, \leq)$ devine o mulțime \emph{total ordonată}.
Totalitatea ordonării provine din proprietatea de \emph{trihotomie}, adică,
date oricare două numere reale $a$ și $b$, exact una dintre $a < b, a > b, a = b$
este adevărată.

Relația de ordine nestrictă are următoarele proprietăți, care sînt relevante
din punctul de vedere al teoriei relațiilor binare:

\begin{itemize}
  \item \emph{reflexibitate}, adică $x \leq x, \forall x \in \RR$;
  \item \emph{antisimetrie}, adică dacă $x \leq y$ și $y \leq x$, rezultă că
    $x$ și $y$ coincid, pentru orice $x, y \in \RR$;
  \item \emph{tranzitivitate}, adică din $x \leq y$ și $y \leq z$ rezultă că
    $x \leq z$, pentru orice $x, y, z \in \RR$.
\end{itemize}

Două alte proprietăți ale mulțimii numerelor reale, care ne arată mai în detaliu
structura sa sînt:

\begin{proposition}[Arhimede]\label{prop:arhimede}
  Fie $x$ un număr real. Există un număr întreg $n$ astfel încît $n \leq x < n+1$.

  Uneori, această proprietate mai este formulată astfel:
  Fie $x < y$ două numere reale. Atunci există $n \in \ZZ$ astfel încît $nx > y$.
\end{proposition}

Cea de-a doua proprietate ne arată legătura cu mulțimea numerelor raționale, pe care
o extinde:

\begin{proposition}[Densitatea lui $\QQ$ în $\RR$]\label{prop:dens-q}
  Fie $x < y$ două numere reale. Atunci există $r \in \QQ$ astfel încît $x < r < y$.

  Cu alte cuvinte, între orice două numere reale există un număr rațional.
\end{proposition}

Pe mulțimea numerelor reale, cel mai adesea vom lucra cu \emph{intervale},
care sînt, din punct de vedere geometric, segmente sau semidrepte luate
din axa numerelor reale.

Un interval (semideschis) se definește astfel:

\[
  [a, b) = \{ x \in \RR \mid a \leq x < b \}.
\]

Majoritatea proprietăților pe care le vom formula pentru mulțimea numerelor
reale vor fi adevărate și pentru orice interval de numere reale, în unele
cazuri fiind necesar să considerăm așa-numita \emph{dreaptă încheiată},
anume $\overline{\RR} = \RR \cup \{ \pm \infty \}$.

Următoarele definiții sînt cunoscute:

\begin{itemize}
  \item $m \in \RR$ se numește \emph{minorant} pentru 
    mulțimea $A \seq \RR$ dacă $m \leq a, \forall a \in A$;
  \item $M \in \RR$ se numește \emph{majorant}  pentru 
    mulțimea $A \seq \RR$ dacă $M \geq a, \forall a \in A$;
  \item Mulțimea $A$ se numește \emph{mărginită} dacă este mărginită și superior,
    și inferior.
  \item $m \in \RR$ se numește \emph{margine inferioară} (infimum) dacă este
    cel mai mare minorant al mulțimii $A$;
  \item $M \in \RR$ se numește \emph{margine superioară} (supremum) dacă este cel
    mai mic majorant al mulțimii $A$.
\end{itemize}

Termenul de \emph{topologie} a dreptei reale, anunțat în titlul secțiunii, ne arată
că putem \qq{să ne mișcăm} pe dreapta reală, cunoscînd ce înseamnă \qq{împrejurimi}.
Formal, ne referim la \emph{vecinătăți}:

\begin{itemize}
  \item Mulțimea $V \seq \RR$ se numește \emph{vecinătate a punctului $x_0 \in \RR$}
    dacă există un interval deschis $I$, astfel încît $x_0 \in I \seq V$;
  \item Mulțimea $V \seq \RR$ se numește \emph{vecinătate a lui $\infty$} dacă există
    un interval deschis $I = (a, \infty) \seq V$;
  \item Mulțimea $V \seq \RR$ se numește \emph{vecinătate a lui $-\infty$} dacă există
    un interval deschis $I = (-\infty, a)$, astfel încît $I \seq V$.
\end{itemize}

Evident, un număr real are o infinitate de vecinătăți: orice interval care îl conține
este o vecinătate a sa. De aceea, pentru un punct $x_0 \in \RR$, vom nota $\kal{V}(x_0)$
mulțimea vecinătăților sale.

Spunem că topologia dreptei reale are o proprietate, numită \emph{separabilitate}, care
ne arată că orice două puncte pot fi separate cu ajutorul vecinătăților. Mai precis,
dacă $x \neq y \in \RR$, atunci există cel puțin o vecinătate a lui $x$ și una a lui $y$
a căror intersecție să fie vidă.

Următoarea noțiune este esențială în studiul analizei matematice:

\begin{definition}\label{def:acumulare}
  Numărul $x_0 \in \overline{\RR}$ se numește \emph{punct de acumulare} al mulțimii $A$
  dacă pentru orice vecinătate $V \in \kal{V}(x_0)$, rezultă $A \cap (V - \{x_0\}) \neq \vid$.

  Un punct $y \in A$ se numește \emph{punct izolat} al mulțimii dacă nu este punct de
  acumulare.
\end{definition}

%%%%%%%%%%%%%%%%%%%%%%%%%%%%%%%%%%%%%%%%%%%%%%%%%%%%%%%%%%%%%%%%%%%%%%

\section{Limite de șiruri}
\label{sec:limite-siruri}

\indent\indent Mai întîi, să ne amintim că prin \emph{șir de numere reale} se înțelege
o funcție $f : \NN \to \RR$. În loc de $f(n)$, vom nota $x_n$ și-l vom numi \emph{al %
  n-lea termen al șirului}. Atunci cînd vrem să ne referim la întreg șirul (i.e.\ la
imaginea funcției $f$), vom nota $(x_n)_{n \in \NN}$ sau, mai simplu, $(x_n)$.

Pentru a studia limita unui șir, avem nevoie de următoarea:
\begin{definition}\label{def:limita-sir}
  Un număr $\ell \in \overline{\RR}$ se numește \emph{limita șirului $(x_n)$} dacă în
  afara unei vecinătăți a lui $\ell$ găsim doar un număr finit de termeni ai șirului.

  Limita unui șir se va nota $\ds\lim_{n \to \infty} x_n$.
\end{definition}

Pentru cazul infinit, mai avem următoarele definiții:
\begin{itemize}
\item $\ds\lim_{n \to \infty} x_n = \infty \Leftrightarrow \forall M \in \RR, \exists %
  N \in \NN \text{ a.î. } x_N > M$. Cu alte cuvinte, oricum am încerca să stabilim
  o margine superioară pentru mulțimea termenilor șirului, sigur găsim cel puțin un
  termen care depășește marginea.
\item $\ds\lim_{n \to \infty} x_n = -\infty \Leftrightarrow \forall m \in \RR, \exists %
  N \in \NN \text{ a.î. } x_N < m$. Cu alte cuvinte, oricum am încerca să găsim o margine
  inferioară pentru mulțimea termenilor șirului, sigur găsim cel puțin un termen care este
  mai mic decît această margine.
\end{itemize}

Evident, dacă există, limita unui șir este unică.

După limită, distingem următoarele tipuri de șiruri:
\begin{itemize}
  \item \emph{șiruri convergente}, care au limită finită;
  \item \emph{șiruri divergente}, care nu au limită sau au limită infinită.
\end{itemize}

Dintr-un șir putem alege \qq{părți}, numite \emph{subșiruri}, iar anumite
proprietăți se păstrează în unele cazuri.

\begin{definition}\label{def:subsir}
  Fie $(x_n)$ un șir de numere reale și $f : \NN^\ast \to \NN^\ast$ o funcție
  strict crescătoare.

  Șirul $\big( n_{f(n)} \big)$ se numește \emph{subșir} al șirului $(x_n)$.
\end{definition}

Cazurile convenabile privitoare la comportamentul subșirurilor sînt:

\begin{itemize}
  \item Orice subșir al unui șir care are limită are aceeași limită cu șirul dat;
  \item Reciproc, dacă toate subșirurile șirului $(x_n)$ au limita $\ell \in \overline{\RR}$,
    atunci șirul $(x_n)$ are aceeași limită;
  \item Fie $(x_n)$ un șir de numere reale cu limita $\ell \in \overline{\RR}$.
    Prin eliminarea sau adăugarea unui număr finit de termeni se obține un șir cu
    aceeași limită.
\end{itemize}

Așa cum vom vedea, ultima proprietate este falsă în cazul seriilor, deoarece în
acea situație, ne vor interesa sumele parțiale ale termenilor șirului, nu
doar termenii.

%%%%%%%%%%%%%%%%%%%%%%%%%%%%%%%%%%%%%%%%%%%%%%%%%%%%%%%%%%%%%%%%%%%%%%

\section{Șiruri convergente}

\id\id Ne vom concentra acum asupra șirurilor convergente, i.e.\ acele
șiruri care au limită finită.

Chiar din definiția convergenței, obținem imediat:

\begin{theorem}\label{thm:conv-marg}
  Orice șir convergent este mărginit.
\end{theorem}

De asemenea, într-un șir convergent, putem lua modulul termenilor, fără
a schimba limita:

\begin{theorem}[Limita modulului]\label{thm:lim-modul}
  Fie $(x_n)$ un șir convergent de numere reale. Atunci șirul $(|x_n|)$ este
  convergent și are loc $\ds\lim_{n \to \infty} |x_n| = | \ds\lim_{n \to \infty} x_n |$.
\end{theorem}

\textbf{Exercițiu:} Ce puteți spune despre reciprocele afirmațiilor de mai sus?

Ne va fi de folos atunci cînd vom formula criterii de convergență următorul
rezultat:

\begin{theorem}[Trecerea la limită în inegalități]\label{thm:lim-ineg}
  Fie $(x_n)$ și $(y_n)$ două șiruri care au limită și astfel încît
  $a_n \leq b_n, \forall n \in \NN$. Atunci $\ds\lim_{n \to \infty} x_n \leq %
  \ds\lim_{n \to \infty} y_n$.
\end{theorem}

Vom mai avea nevoie și de:
\begin{definition}\label{def:sir-monoton}
  Șirul $(x_n)$ se numește \emph{monoton crescător} dacă $x_n \leq x_{n + 1}$,
  pentru orice $n \in \NN$. Echivalent, $\frac{x_n}{x_{n+1}} \leq 1$.

  Șirul se numește \emph{monoton descrescător} dacă $x_n \geq x_{n + 1}$, pentru
  orice $n \in \NN$. Echivalent, $\frac{x_n}{x_{n+1}} \geq 1$.

  Dacă un șir este fie crescător, fie descrescător, el se numește \emph{monoton}.
\end{definition}

În particular, obținem:

\begin{itemize}
  \item Dacă $(x_n)$ este un șir crescător și convergent, iar $\ell$ este
    limita sa, atunci $x_n \leq \ell, \forall n \in \NN$;
  \item Similar, dacă $(x_n)$ este descrescător și convergent, iar $\ell$
    este limita sa, atunci $x_n \geq \ell, \forall n \in \NN$.
\end{itemize}

Foarte util în exerciții este următorul rezultat:

\begin{theorem}[K. Weierstrass]\label{thm:mon-marg-conv}
  Dacă un șir este monoton și mărginit, atunci el este convergent.
\end{theorem}

Așadar, pentru a verifica convergența, vom verifica monotonia și mărginirea.

%%%%%%%%%%%%%%%%%%%%%%%%%%%%%%%%%%%%%%%%%%%%%%%%%%%%%%%%%%%%%%%%%%%%%%

\section{Criterii de existență a limitei și de convergență}

\id\id Cel mai clar criteriu de convergență, dar și cel mai impractic
de folosit în exerciții este următorul:

\begin{theorem}[Criteriul cu $\varepsilon$]\label{thm:criteriu-epsilon}
  Fie $(x_n)$ un șir de numere reale. Un număr $\ell \in \RR$ este limita șirului
  $(x_n)$ dacă și numai dacă pentru orice $\varepsilon > 0$, există un rang $N_\varepsilon \in \NN^\ast$,
  astfel încît $|x_n - \ell| < \varepsilon$, pentru orice $n \geq N_\varepsilon$.
\end{theorem}

Acest criteriu poate fi adaptat și pentru limite infinite:

\begin{theorem}[Criteriul cu $\varepsilon$ pentru limite infinite]\label{thm:criteriu-eps-inf}
  Fie $(x_n)$ un șir de numere reale.
  \begin{itemize}
    \item $\ds\lim_{n \to \infty} x_n = \infty$ dacă și numai dacă pentru orice $\varepsilon > 0$, există
      un rang $N_\varepsilon \in \NN^\ast$, astfel încît $x_n > \varepsilon$, pentru orice $n \geq N_\varepsilon$;
    \item $\ds\lim_{n \to \infty} x_n = -\infty$ dacă și numai dacă pentru orice $\varepsilon > 0$,
      există un rang $N_\varepsilon \in \NN^\ast$, astfel încît $x_n < -\varepsilon$, pentru orice $n \geq N_\varepsilon$.
  \end{itemize}
\end{theorem}

Există o noțiune ceva mai relaxată de convergență, numită \emph{convergență în sens Cauchy}:

\begin{definition}\label{def:cauchy}
  Fie $(x_n)$ un șir de numere reale. Spunem că el este \emph{convergent în sens Cauchy} (sau că are %
  proprietatea Cauchy) dacă termenii săi pot fi făcuți oricît de apropiați, de la un rang încolo.

  Formal, $\forall \varepsilon, \exists N_\varepsilon$ astfel încît pentru orice $m, n \geq N_\varepsilon$,
  să avem $|x_m - x_n| < \varepsilon$.
\end{definition}

\textbf{Exercițiu:} Găsiți o legătură între convergența în sens Cauchy și convergența \qq{normală}.

Operațiile permise cu \emph{șirurile convergente} sînt următoarele:
\begin{itemize}
  \item Limita sumei este egală cu suma limitelor;
  \item Limita produsului este egală cu produsul limitelor;
  \item Limita cîtului este egală cu cîtul limitelor (cînd cîtul are sens);
  \item Limita unei exponențiale se distribuie:
    \[
      \lim_{n \to \infty} (x_n)^{y_n} = \Big( \lim_{n \to \infty} x_n \Big)^{\lim_{n \to \infty} y_n}
    \]
\end{itemize}

Atunci cînd șirurile nu sînt convergente, trebuie ținut cont de următoarele
\emph{cazuri de nedeterminare}: $\frac{0}{0}, \frac{\pm\infty}{\pm\infty} 1^\infty, \infty - \infty, \pm\infty \cdot 0, %
0^0, \infty^0$.


Ajungem, în fine, la criterii de convergență.

\begin{itemize}
  \item \textbf{Criteriul comparației:} Dacă un șir este mai mic (termen cu termen) decît un șir
    convergent, atunci și acest șir este convergent. Similar, dacă este mai mare (termen cu termen)
    decît un șir divergent spre $\infty$, este divergent (respectiv mai mic decît un șir divergent
    spre $-\infty$).
  \item Fie $(x_n), (y_n)$ două șiruri de numere reale, astfel încît $x_n \to 0$ și $y_n$ este mărginit.
    Atunci $x_n \cdot y_n \to 0$.
  \item \textbf{Criteriul cleștelui}: Fie $(x_n), (y_n), (z_n)$ șiruri de numere reale, astfel încît
    $x_n \leq y_n \leq z_n$, pentru orice $n \in \NN^\ast$. Dacă $x_n, z_n \to \ell$, atunci $y_n \to \ell$.
  \item \textbf{Limite remarcabile:}
    \[
      \lim_{n \to \infty} \frac{\sin x_n}{x_n} = \lim_{n \to \infty} \frac{\tan x_n}{x_n} = %
      \lim_{n \to \infty} \frac{\arcsin x_n}{x_n} = \lim_{n \to \infty} \frac{\arctan x_n}{x_n} = 1.
    \]
  \item \textbf{Stolz-Cesaro:} Fie $(x_n)$ și $(y_n)$ două șiruri de numere reale, cu proprietățile:
    \begin{itemize}
      \item $(y_n)$ strict crescător și nemărginit, cu termeni nenuli;
      \item Șirul $(z_n)$, definit prin $z_n = \frac{x_{n + 1} - x_n}{y_{n + 1} - y_n}$ are limita
        $\ell \in \overline{\RR}$.
    \end{itemize}
     Atunci șirul $\Big( \frac{x_n}{y_n} \Big)$ are limita egală cu $\ell$.
     
     Rezultatul se păstrează și dacă $\lim x_n = \lim y_n = 0$.
   \item \textbf{Criteriul raportului (D'Alembert)}: Fie $(x_n)$ un șir de numere reale nenule și
     $L = \ds\lim_{n \to \infty} \Big| \frac{x_{n + 1}}{x_n} \Big|$. Atunci:
     \begin{itemize}
       \item Dacă $L < 1$, șirul este convergent;
       \item Dacă $L > 1$, șirul este divergent;
       \item Dacă $L = 1$, criteriul nu decide.
     \end{itemize}
   \item \textbf{Criteriul radical (Cauchy)}: Fie $(x_n)$ un șir de numere reale și
     $L = \ds\lim_{n \to \infty} \sqrt[n]{|x_n|}$. Atunci:
     \begin{itemize}
       \item Dacă $L < 1$, șirul este convergent;
       \item Dacă $L > 1$, șirul este divergent;
       \item Dacă $L = 1$, criteriul nu decide.
     \end{itemize}
\end{itemize}

\begin{definition}[Numărul $e$]\label{def:e}
  Fie $(x_n)$ un șir de numere reale, cu $\ds\lim_{n \to \infty} x_n = \infty$. Atunci:
  \[
    \lim_{n \to \infty} \Big( 1 + \frac{1}{x_n} \Big)^{x_n} = e.
  \]
  
  Similar, dacă $(y_n)$ este un șir de numere reale, cu $\ds\lim_{n \to \infty} y_n = 0$,
  atunci:
  \[
    \lim_{n \to \infty} \Big( 1 + y_n \Big)^{\frac{1}{y_n}} = e.
  \]
\end{definition}

Practic, numărul $e$ (constanta lui Euler) vine să rezolve nedeterminarea $1^\infty$.

%%%%%%%%%%%%%%%%%%%%%%%%%%%%%%%%%%%%%%%%%%%%%%%%%%%%%%%%%%%%%%%%%%%%%%

\section{Exerciții}

1. Precizați valoarea de adevăr a propozițiilor (în caz de falsitate, dați un contraexemplu):
\begin{enumerate}[(a)]
  \item Orice șir monoton este mărginit.
  \item Orice șir mărginit este monoton.
  \item Orice șir convergent este monoton.
  \item Orice subșir al unui șir monoton este monoton.
  \item Există șiruri monoton crescătoare care au cel puțin un subșir monoton descrescător.
  \item Suma a două șiruri monotone este un șir monoton.
  \item Diferența a două șiruri monoton crescătoare este un șir monoton crescătoare este un șir
    monoton crescător.
  \item Suma a două șiruri nemărginite este un șir nemărginit.
  \item Orice șir divergent este nemărginit.
  \item Orice șir nemărginit are limita $\pm \infty$.
  \item Produsul a două șiruri care nu au limită este un șir care nu are limită.
  \item Dacă două șiruri sînt convergente, atunci raportul lor este un șir convergent.
  \item Dacă pătratul unui șir $(x_n)$ este convergent, atunci șirul $(x_n)$ este convergent.
  \item Dacă un șir convergent are toți termenii nenuli, atunci limita sa este nenulă.
\end{enumerate}

2. Să se studieze convergența și să se calculeze limitele șirurilor:
\begin{multicols}{2}
\begin{enumerate}[(a)]
  \item $x_n = n^2 - n$;
  \item $x_n = -n^3 + 2n^2$;
  \item $x_n = -4n^3 + 3n - 1$;
  \item $x_n = 4n^4 - 5n^6 + 3$;
  \item $x_n = \frac{n}{n^2 + 1}$;
  \item $x_n = \frac{2n^2 - 3n}{4n^2 - n + 1}$;
  \item $x_n = \frac{2n^3 + 5n - 1}{3n^2 + n + 1}$;
  \item $x_n = \frac{7 - 2n - n^3}{2n^2 + 3n + 1}$;
  \item $x_n = \Big(\cos \frac{1}{n} \Big)^n$;
  \item $x_n = (n^2 + 1)^{\frac{n}{n^2 + 1}}$;
  \item $x_n = \Big( \frac{n}{n^2 + 3} \Big)^{\frac{2}{n}}$;
  \item $x_n = \frac{\ln n}{n}$;
  \item $x_n = \frac{1}{n} \Big(1 + \frac{1}{2} + \frac{1}{3} + \dots + \frac{1}{n} \Big)$;
  \item $x_n = \frac{1^4 + 2^4 + 3^4 + \dots + n^4}{n^5}$;
  \item $x_n = \frac{1}{n} \Big( 1 + \frac{1}{3} + \frac{1}{5} + \dots + \frac{1}{2n - 1} \Big)$;
  \item $x_n = \frac{(-1)^n \cdot n}{n + 2}$;
  \item $x_n = \frac{1}{1 \cdot 4} + \frac{1}{2 \cdot 5} + \dots + \frac{1}{n(n + 3)}$;
  \item $x_n = \Big(1 + \frac{2}{n} \Big)^n$;
  \item $x_n = \Big(1 + \frac{1}{n^2 + 1} \Big)^{n^2}$;
  \item $x_n = \Big(\frac{n + 5}{n + 2} \Big)^n$;
  \item $x_n = \Big( \frac{n^4 + 2n^2 + 1}{n^4 + n^2 + 3n} \Big)^{\frac{n^2}{n + 1}}$;
  \item $x_n = \frac{4n^2 + 1}{n^2 + 1}$;
  \item $x_n = \frac{1}{n + 1} + \frac{1}{n + 2} + \frac{1}{n + 3}$;
  \item $x_n = \frac{1}{n} \sin(n^2)$;
  \item $x_n = \frac{n \cdot \sin n}{n^2 + 1}$;
  \item $x_n = \frac{1}{2} + \frac{1}{3} + \frac{1}{4} + \frac{1}{5} + \dots + \frac{1}{n}$.
\end{enumerate}
\end{multicols}


\end{document}
